In previous sections we intuitively introduced the concept of a confidence interval and evaluated basic interval estimations for the population mean. However, in real statistical analysis and data science projects, we rarely limit ourselves to studying a single population or assuming known population parameters. We frequently need to compare two treatments (such as in A/B testing for websites), evaluate proportions or conversion rates, analyze data variability or dispersion, and, critically before any study begins, accurately calculate the sample size required to achieve a desired statistical precision or power.

In this section we formally, thoroughly, and rigorously develop the theory of interval estimation for differences in means, standard errors, proportions, variances, and sample size design.

\section{Estimation of the Mean and Difference Between Two Means}

When two populations are studied, primary interest often centers on estimating the difference between their means, $\mu_1 - \mu_2$, using independent or paired samples drawn from each population. The natural, minimum-variance point estimator for this difference is the difference between sample means: $\bar{X}_1 - \bar{X}_2$.

The construction of the confidence interval for $\mu_1 - \mu_2$ depends critically on the sampling design and the assumptions or knowledge regarding the population variances ($\sigma_1^2, \sigma_2^2$).

\begin{teorema}[Case 1: Independent Samples with Known Variances]
	Let $X_{11}, X_{12}, \ldots, X_{1n_1}$ and $X_{21}, X_{22}, \ldots, X_{2n_2}$ be two independent random samples of sizes $n_1$ and $n_2$ drawn from populations with means $\mu_1, \mu_2$ and known variances $\sigma_1^2, \sigma_2^2$. If the populations are normal or the sample sizes are large ($n_1, n_2 \geq 30$ by the Central Limit Theorem), a $(1-\alpha)100\%$ confidence interval for the difference of means is:
	\begin{align}
		(\bar{X}_1 - \bar{X}_2) - Z_{\alpha/2} \sqrt{\frac{\sigma_1^2}{n_1} + \frac{\sigma_2^2}{n_2}} \leq \mu_1 - \mu_2 \leq (\bar{X}_1 - \bar{X}_2) + Z_{\alpha/2} \sqrt{\frac{\sigma_1^2}{n_1} + \frac{\sigma_2^2}{n_2}},
	\end{align}
	where $Z_{\alpha/2}$ is the critical value from the standard normal distribution leaving an upper tail area of $\alpha/2$.
\end{teorema}

\begin{teorema}[Case 2: Independent Samples with Unknown but Equal Variances]
	If the population variances are unknown but homoscedasticity can be assumed ($\sigma_1^2 = \sigma_2^2 = \sigma^2$), we pool the information from both samples to obtain a degrees-of-freedom-weighted estimate of the common variance, denoted by the pooled variance $S_p^2$:
	\begin{align}
		S_p^2 = \frac{(n_1 - 1)S_1^2 + (n_2 - 1)S_2^2}{n_1 + n_2 - 2},
		\label{eq:varianza_combinada}
	\end{align}
	where $S_1^2$ and $S_2^2$ are the respective unbiased sample variances. 
	
	The standardized statistic follows a Student's $t$-distribution with $\nu = n_1 + n_2 - 2$ degrees of freedom. The $(1-\alpha)100\%$ confidence interval for $\mu_1 - \mu_2$ is:
	\begin{align}
		(\bar{X}_1 - \bar{X}_2) \pm t_{\alpha/2, \, n_1+n_2-2} \, S_p \sqrt{\frac{1}{n_1} + \frac{1}{n_2}}.
	\end{align}
\end{teorema}

\begin{teorema}[Case 3: Independent Samples with Unknown and Unequal Variances (Welch)]
	When population variances are unknown and it is not reasonable to assume they are equal ($\sigma_1^2 \neq \sigma_2^2$), we cannot use the pooled variance. In this scenario (known as the Behrens-Fisher problem), we use the individual sample variances and approximate the distribution of the test statistic using Student's $t$-distribution with Welch-Satterthwaite effective degrees of freedom:
	\begin{align}
		\nu = \frac{\left( \frac{S_1^2}{n_1} + \frac{S_2^2}{n_2} \right)^2}{\frac{\left(S_1^2 / n_1\right)^2}{n_1 - 1} + \frac{\left(S_2^2 / n_2\right)^2}{n_2 - 1}}.
		\label{eq:df_welch}
	\end{align}
	If $\nu$ is not an integer, it is typically rounded down to the nearest integer to maintain a conservative attitude regarding interval coverage. The $(1-\alpha)100\%$ interval is given by:
	\begin{align}
		(\bar{X}_1 - \bar{X}_2) \pm t_{\alpha/2, \, \nu} \sqrt{\frac{S_1^2}{n_1} + \frac{S_2^2}{n_2}}.
	\end{align}
\end{teorema}

\begin{definicion}[Case 4: Paired or Dependent Sampling]
	When two samples are pairwise correlated (for example, measurements from the same patient or system before and after a treatment), we cannot treat them as independent. Instead, we reduce the two-variable problem to a single-variable problem by computing individual paired differences:
	\begin{align}
		D_i = X_{1i} - X_{2i}, \quad \text{for } i = 1, 2, \ldots, n.
	\end{align}
	
	We compute the sample mean of the differences $\bar{D} = \frac{1}{n}\sum D_i$ and the sample standard deviation of the differences $S_D = \sqrt{\frac{1}{n-1}\sum (D_i - \bar{D})^2}$. The $(1-\alpha)100\%$ confidence interval for the population mean difference $\mu_D = \mu_1 - \mu_2$ is:
	\begin{align}
		\bar{D} - t_{\alpha/2, \, n-1} \frac{S_D}{\sqrt{n}} \leq \mu_D \leq \bar{D} + t_{\alpha/2, \, n-1} \frac{S_D}{\sqrt{n}}.
	\end{align}
\end{definicion}
